% Chapter 6, typeset from lectures/L06: the README is the opener and the closing review, and each
% part of Appendix A is a section, in the same order and with the same subsections. Appendix B,
% the exercises, is the last section.

\chapter{Counters and Shift Registers}
\label{c6:ch}

{\large\itshape Building blocks for anything that needs to count or move bits.\par}

\begin{chapterbox}{In this chapter}
\begin{itemize}
  \item Counters: a register, an adder and a constant \code{1}, and why wraparound comes for free.
  \item Shift registers: moving data one bit per clock edge, in one direction used consistently.
  \item Serial-in/parallel-out and parallel-in/serial-out, and which fits which job.
  \item An 8-bit serial receiver: a shift register plus a bit counter, knowing when a byte has
    arrived.
  \item Two demonstrations on the DE0-CV: a row of LEDs counting and wrapping, and a byte shifted in
    one button press at a time.
\end{itemize}
\end{chapterbox}

This chapter builds four things, in this order. A 4-bit counter in CircuitVerse: a register, an
adder, and a constant \code{1}, run past \code{1111} to watch it wrap. The same counter in VHDL. A
shift register in VHDL, traced one clock edge at a time on paper before any of it is written. And
\code{serial_rx8}, the shift register just written plus the bit counter that tells it when a byte
is done: the two halves of the chapter meeting in one module. Then both go onto the board. First
the counter, its top bits on a row of LEDs, slowed to something the eye can follow, running up to
all-ones and wrapping back to zero on its own; then \code{serial_rx8} with \code{shift_enable}
driven by a push button instead of held high, so one press shifts one bit and the byte assembles
across eight LEDs, most significant bit first.

Before the counter in \secref{c6:a1} reaches \code{1111}, predict what it will show next, and then
find the element in the drawing that makes it happen. There isn't one, which is the point: nothing
in the circuit performs the wrap, and that is a consequence of the register's width rather than
logic anybody added.

\code{serial_rx8} (\file{lectures/L06/serial_rx8/serial_rx8.vhd}) is the first design in the book
whose correctness you cannot check by eye at full speed: enabled continuously it swallows a byte in
160\,ns. That is why it ships a reference testbench, and why the board demonstration drives
\code{shift_enable} from a button instead. Slowing the \emph{clock} would be the software instinct
and the wrong move; the module already has the control that lets you take one bit at a time. The
board demonstration in \secref{c6:a7} wraps it in Chapter~\ref{c4:ch}'s \code{reset_sync} and
\code{button_sync}, both unchanged, so it is a composition exercise as much as a receiver. Read the
sections through before you start on the exercises; running their testbenches is covered in
\secref{c2:app:c}.

% Sections 6.1 to 6.7, from lectures/L06/appendix/a_counters_and_shift_registers.md: one section
% per part, A.1 to A.7.

\section{From Registers to Counters}
\label{c6:a1}
A register (Chapter~\ref{c3:ch}) is a bank of D flip-flops storing whatever value is fed in on the
next rising edge. A \textbf{counter} is the same idea with one twist: instead of a new value from
outside, you feed the register its \emph{own current value plus one}. Each edge it captures
\code{current_value + 1}, so the stored number increments once per clock cycle with no further
input.

\subsection{Build it by hand first}
\label{c6:byhand}
As with every circuit in this book, build it before you write it. A 4-bit counter is three
elements, and watching it run is the point:
\begin{itemize}
  \item A \textbf{4-bit register}: four D flip-flops sharing one clock, exactly the register from
    \secref{c3:a5}.
  \item An \textbf{adder} and a \textbf{constant \code{1}}: the register's four outputs into one
    adder input, the constant into the other.
  \item The adder's sum wired back to the register's four \code{D} inputs. That feedback path is
    the whole circuit: the register holds a number, the adder computes that number plus one, the
    next edge stores it. It is \secref{c3:a1}'s feedback loop again, with an adder in the path
    instead of a plain wire.
\end{itemize}

\bookfigure{counter_circuit}{A 4-bit counter: a register whose own output, plus one, feeds back to
its input. The adder's \code{carry_out} is connected to nothing.}{c6:fig:counter}

\noindent Then confirm, one clock edge at a time, with a period of \code{1000 ms}:
\begin{itemize}
  \item Watch the outputs step \code{0000, 0001, 0010, 0011, ...}, one count per rising edge.
  \item \textbf{Let it run past \code{1111}.} It returns to \code{0000} on its own, and nothing in
    your circuit says to do that: there is no comparison, no reset logic, no ``if the count is at
    maximum'' anywhere. The adder produced \code{10000}, the fifth bit had nowhere to go, and the
    four bits that survived are \code{0000}. That is \textbf{overflow}, and in a counter it is not
    a bug to work around but the feature the rest of this chapter and all of Chapter~\ref{c7:ch}
    are built on.
  \item \textbf{Watch \code{carry_out} on the adder while it happens.} That is the fifth bit, and
    the drawing shows exactly where it ends up: high throughout the single count at which the
    register rolls over, connected to nothing, leaving the adder and stopping there. A wider
    counter would feed it into a fifth stage; here it is discarded, and discarding it \emph{is} the
    wraparound. So ``the count wraps because it runs out of bits'' is not a figure of speech. You
    can point at the bit it ran out of.
\end{itemize}
Save the project. Chapter~\ref{c7:ch} starts from this exact circuit and adds two things to it,
and reopening it beats redrawing it.

\subsection{The same counter in VHDL}
\label{c6:vhdlcounter}
A signal and a process, exactly like Chapter~\ref{c3:ch}'s registers, except the right-hand side
refers to the signal itself:

\begin{vhdlcode}
signal counter: natural range 0 to 15;
...
process(clock, reset_s2_n) is
begin
    if (reset_s2_n = '0') then
        counter <= 0;
    elsif (rising_edge(clock)) then
        counter <= counter + 1;
    end if;
end process;
\end{vhdlcode}

\noindent\code{natural} is a predefined VHDL integer subtype: a non-negative whole number, as
opposed to the \code{std_logic} types that model bits directly. VHDL's arithmetic operators work on
it the way ordinary integer maths does, which is why \code{counter + 1} needs no extra machinery.
Synthesis still turns it into a register of flip-flops, exactly as a \code{std_logic_vector} would;
the difference is purely how you write and reason about the value.

The range constraint \code{0 to 15} is not just documentation. It tells the synthesis tool this
counter needs four bits, so incrementing past \code{15} rolls back to \code{0} for free, purely as a
consequence of running out of bits, with no explicit ``if counter = 15'' logic in the drawing or the
VHDL.

That free wrap has a precondition worth stating: it only happens when the upper bound is
\code{2^N - 1}. \code{0 to 15} wraps for free; \code{0 to 9} does not, because \code{10} is not a
power-of-two boundary, and a modulo-10 counter needs an explicit comparison. That is exactly what
Exercise~\ref{c6:ex:4} asks you to write.

\textbf{One place the simulator disagrees with the hardware.} The free wrap is a property of the
\emph{synthesized} circuit, and the drawing shows it is real. A VHDL \emph{simulator} is stricter
than a wire:
\begin{itemize}
  \item \code{natural range 0 to 15} is range-checked, so GHDL treats \code{15 + 1 = 16} as an
    error and aborts rather than wrapping.
  \item The synthesized hardware has no such notion. There is no sixteenth value for the bits to
    take, so they land on \code{0000}, exactly as your four LEDs did.
\end{itemize}
The gate network and the chip both wrap; only the simulator objects, because a range-constrained
integer is a promise you made and \code{16} breaks it. The rule that follows is not ``trust the
hardware'': it is that anything you intend to verify should say what it means. Write the comparison
out, the way \file{timer.vhd} and \file{serial_rx8.vhd} do, and the drawing, the simulator and the
chip all agree.

This wraparound-by-bit-width is the single most important property of a counter, and it is what
Chapter~\ref{c7:ch} builds a timer out of: a counter that watches for one particular count and
raises a flag when it arrives.

\section{Shift Registers}
\label{c6:a2}
A \textbf{shift register} is, again, N D flip-flops, wired so each flip-flop's output feeds the
\emph{next} one's input rather than its own. On every clock edge every bit moves one position down
the chain simultaneously:

\begin{console}
serial_in -> [D Q]-->[D Q]-->[D Q]-->[D Q] -> serial_out
              FF0     FF1     FF2     FF3
             bit 0   bit 1   bit 2   bit 3
               ^clock  ^clock  ^clock  ^clock  (all share the same clock)
\end{console}

\noindent After one edge, \code{FF0}'s contents are in \code{FF1}, \code{FF1}'s in \code{FF2}, and
so on; a new bit enters at \code{FF0} from \code{serial_in}, and whatever was in the last flip-flop
falls off the end as \code{serial_out}. Run it for N edges and N serial bits have moved all the way
through.

\textbf{One direction, used everywhere in this book.} A shift register can move bits either way,
and both are equally valid; mixing them within one design is what causes confusion. So every shift
register here follows the diagram above:
\begin{itemize}
  \item the chain maps onto a \code{std_logic_vector} with \code{FF0} as \textbf{bit 0} and
    \code{FF(N-1)} as bit \code{N-1}.
  \item serial data \textbf{enters at bit 0} and \textbf{leaves from bit \code{N-1}}.
  \item every bit moves one position toward the \textbf{most significant} bit per edge.
\end{itemize}
In VHDL that direction is a single expression, identical in every example below:

\begin{vhdlcode}
shift_reg <= shift_reg(N-2 downto 0) & serial_in;
\end{vhdlcode}

\noindent\code{shift_reg(N-2 downto 0)} is everything except the top bit, moved up one position;
the top bit has nowhere to go and falls off as \code{serial_out}; \code{serial_in} fills the vacated
bit 0.

Whenever you meet a shift register outside this book, check its direction before assuming. The
mirror-image form, \code{serial_in & shift_reg(N-1 downto 1)}, is just as common and reads almost
identically at a glance.

Shift registers are named by how data enters and leaves:

\begin{booktable}{@{}>{\raggedright\arraybackslash}p{7.6em}
  >{\raggedright\arraybackslash}p{6.2em}>{\raggedright\arraybackslash}p{6.2em}L@{}}
  \thead{Configuration} & \thead{Data in} & \thead{Data out} & \thead{Typical use} \\
  \midrule
  SISO (serial-in/\allowbreak serial-out) & one bit per clock & one bit per clock & a pure delay
    line \\
  \textbf{SIPO} (serial-in/\allowbreak parallel-out) & one bit per clock & all N bits at once &
    receiving a stream of serial bits, e.g.\ an incoming byte on a single wire \\
  \textbf{PISO} (parallel-in/\allowbreak serial-out) & all N bits at once (loaded) & one bit per
    clock & sending N bits out over a single wire, e.g.\ driving a chain of LEDs or an external
    shift-register chip \\
  PIPO (parallel-in/\allowbreak parallel-out) & all N bits at once & all N bits at once &
    effectively a plain register (\secref{c3:a5}) \\
\end{booktable}

\noindent This chapter covers the two practically useful ones, SIPO and PISO.

\section{Serial-In/Parallel-Out (SIPO)}
\label{c6:a3}
A SIPO register is what you reach for when bits arrive one at a time on a single wire and you need
the accumulated value all at once. The classic case is deserializing an incoming byte, which is
what the receive path of any serial protocol does at heart.

Which end of the byte arrives first is each protocol's own decision. SPI and CAN send the most
significant bit first by default, which is what \secref{c6:a2}'s direction assumes: the first bit
to arrive ends up in the MSB once all N have shifted in. (SPI bit order is configurable on most
controllers, so check the part.) A protocol sending the least significant bit first, such as UART,
shifts the mirror-image way, and its idiom is the mirror image too.

The idiom is one line inside a clocked process, \secref{c6:a2}'s shift expression verbatim:

\begin{vhdlcode}
signal shift_reg: std_logic_vector(7 downto 0);
...
process(clock, reset_s2_n) is
begin
    if (reset_s2_n = '0') then
        shift_reg <= (others => '0');
    elsif (rising_edge(clock)) then
        shift_reg <= shift_reg(6 downto 0) & serial_in;
    end if;
end process;
parallel_out <= shift_reg;
\end{vhdlcode}

\noindent After 8 edges the received bits have fully replaced the register's original contents,
and \code{parallel_out} exposes all 8 at once.

\section{Parallel-In/Serial-Out (PISO)}
\label{c6:a4}
A PISO register reverses SIPO's job: load N bits at once, then send them out one at a time. It is
what you reach for when a fixed pattern has to go out over a single wire, such as a chain of LEDs
or a display driven by an external shift-register chip.

It needs one more control signal than SIPO: something to distinguish \emph{loading} a new parallel
value from \emph{shifting} the existing one out.

\begin{vhdlcode}
process(clock, reset_s2_n) is
begin
    if (reset_s2_n = '0') then
        shift_reg <= (others => '0');
    elsif (rising_edge(clock)) then
        if (load = '1') then
            shift_reg <= parallel_in;                          -- Parallel load.
        elsif (shift = '1') then
            shift_reg <= shift_reg(N-2 downto 0) & serial_in;  -- Shift toward the MSB.
        end if;
    end if;
end process;
serial_out <= shift_reg(N-1);
\end{vhdlcode}

\noindent Bit \code{N-1} is always the next bit to leave, so a loaded value goes out
most-significant bit first: load \code{"1010"} into a 4-bit register and \code{serial_out} presents
\code{1}, \code{0}, \code{1}, \code{0} over the following cycles.

Notice how little separates this from \secref{c6:a3}'s SIPO. The \textbf{shift expression is
identical}, same direction, same line of VHDL. Only two things differ: PISO adds a \code{load}
branch, and it taps \code{shift_reg(N-1)} as a serial output where SIPO exposes the whole vector.
``SIPO'' and ``PISO'' name how you \emph{wire up} a shift register, not two kinds of hardware,
which is why \secref{c6:a2} could describe the shift once before either name appeared.

Whether a design needs a genuinely new \code{serial_in} on every shift, or is happy to recirculate
whatever falls off the end, is up to how it is wired. That recirculation is exactly the trick the
walking-LED example in Chapter~\ref{c7:ch} uses, once it has a timer to pace the shifting.

\section{Worked Example: An 8-Bit Serial Receiver}
\label{c6:a5}
The two halves of this chapter meet in one module. A shift register can accept serial bits forever
but cannot tell you \emph{when} a complete byte has arrived, because it has no notion of ``how many
bits so far''. That is a counting problem, and \secref{c6:a1} already solved it.
\file{lectures/L06/serial_rx8/serial_rx8.vhd} combines the two, plus a one-cycle
\code{data_ready} pulse when the eighth bit lands.

This is not a toy: it is the front half of every serial receiver there is. SPI, CAN and UART receive
paths all start with a shift register and a bit counter, differing only in what decides when a bit
is valid.

\begin{booktable}{@{}p{6.4em}p{4.8em}p{9.6em}L@{}}
  \thead{Port} & \thead{Direction} & \thead{Type} & \thead{Description} \\
  \midrule
  \code{clock} & in & \code{std_logic} & System clock. \\
  \code{reset_s2_n} & in & \code{std_logic} & Active-low, \textbf{already-synchronized} reset;
    clears the register and the bit count. Asynchronous, so it clears them with no clock edge
    involved. \\
  \code{shift_enable} & in & \code{std_logic} & High for exactly one clock cycle per incoming
    bit. \\
  \code{serial_in} & in & \code{std_logic} & The incoming bit, sampled on each enabled rising
    edge. \\
  \code{data_out} & out & \code{std_logic_vector(7 downto 0)} & The byte received so far, most
    significant bit first. \\
  \code{data_ready} & out & \code{std_logic} & One-cycle pulse as the eighth bit lands. \\
\end{booktable}

\textbf{The reset port is the synchronized one, and the \code{_s2} says so.} It is named
\code{reset_s2_n} rather than \code{reset_n} because what belongs on it is the output of a
\code{reset_sync} (\secref{c4:a5}), the module you wrote in Chapter~4's Exercise~\ref{c4:ex:8}:
assert asynchronously, release synchronously. \textbf{Instantiate one in whatever design encloses
this receiver and wire its \code{reset_s2_n} here}, which is exactly what \secref{c6:a7}'s board
wrapper does. Wire a push button straight to this port instead and you get the race
Chapter~\ref{c4:ch} exists to prevent: the shift register and the bit counter are a dozen
flip-flops, and an asynchronously \emph{released} reset lets them leave reset on different edges,
so the receiver can start counting a byte a cycle before it starts shifting one. The rule is the
same one \secref{c7:a3} states for \code{timer}, and it holds for every subcomponent in this book:
a module composed inside a larger design never touches a raw asynchronous input itself.

\code{counter}, \code{sipo8} and \code{piso8} in \secref{c6:sec:exercises} take \code{reset_s2_n}
for the same reason. None of them is a design you put on a board by itself: a counter, a shift
register and a receiver are parts, and a part is entitled to assume the design around it has done
the synchronizing. So the whole of this chapter takes the synchronized reset, and only the wrappers
in \secref{c6:a6} and \secref{c6:a7}, which own the pins, ever see a raw \code{reset_n}.

A \textbf{top level} is named the other way round: it takes \code{reset_n} and instantiates
\code{reset_sync} itself, because it is the thing holding the pin. \code{blinker} and
\code{walking_led} in Chapter~\ref{c7:ch}, and both Chapter~\ref{c8:ch} capstones, are written that
way.

Get into the habit of reading the \code{_s2} as a requirement on the caller, because from here on
most of what you write is a subcomponent.

On each rising edge where \code{shift_enable = '1'}, \code{serial_in} is shifted into bit 0 by
\secref{c6:a3}'s idiom unchanged, so the first bit to arrive has walked up to bit 7 by the time the
eighth lands, and an internal \code{natural range 0 to 7} bit counter (\secref{c6:a1}) advances. On
the edge where that counter is already at \code{7}, \code{data_ready} pulses high for one cycle and
the counter is cleared back to \code{0}, ready for the next byte. \code{data_out} is driven
concurrently from the shift register, so the complete byte is on it during the cycle
\code{data_ready} is high.

Note that the counter is \emph{cleared} here rather than left to wrap: \code{0 to 7} is a
power-of-two range so it would wrap for free, but writing it out keeps the simulator in agreement
with the hardware, per \secref{c6:a1}.

Two details are worth dwelling on, because both are patterns you have met more than once:
\begin{itemize}
  \item \textbf{\code{data_ready} is a pulse, not a level.} It is driven low unconditionally at the
    top of the clocked branch and raised only on the edge completing a byte: the same single-cycle
    shape as \secref{c3:a7}'s edge detector. A consumer that misses it has missed the byte, which
    is exactly why it is paired with a \code{data_out} that \emph{holds} its value afterwards. A
    real peripheral would put a holding register or a small FIFO here so a slow consumer cannot
    lose data.
  \item \textbf{\code{shift_enable} gates the counter, not just the shift.} Both advance together
    or neither does, so a gap in the incoming stream pauses the receiver mid-byte rather than
    corrupting it, and it resumes on exactly the bit it was waiting for.
\end{itemize}
This is the whole of \file{serial_rx8.vhd}:

\vhdlfile{lectures/L06/serial_rx8/serial_rx8.vhd}

\textbf{Checking it:} the example ships \file{serial_rx8_tb.vhd}, which shifts two bytes in with a
pause between them and checks that \code{data_ready} pulses once per byte, never early:

\begin{shell}
cd lectures/L06/serial_rx8
ghdl -a --std=93 serial_rx8.vhd serial_rx8_tb.vhd
ghdl -e --std=93 serial_rx8_tb
ghdl -r --std=93 serial_rx8_tb --assert-level=error --stop-time=10ms
\end{shell}

\noindent Clocked at \code{50 MHz} and left enabled, this receiver consumes a byte in 160\,ns,
which no LED can show you. That is why the testbench matters here in a way it did not for a
blinking LED, and it is the case for most real digital design. But it is not a reason to skip the
board, because \code{shift_enable} is exactly the way out: see \secref{c6:a7}.

\section{FPGA Demonstration: Watching a Counter Wrap}
\label{c6:a6}
The course demonstrates the counter on the DE0-CV, built the same way as every design since
Chapter~\ref{c1:ch}: a Quartus Prime Lite project targeting \code{5CEBA4F23C7N}, the \file{.vhd}
file added, compile, assign each port to a physical pin in the Pin Planner, then recompile and
program the board (Appendix~\ref{app:quartus}).

One thing has to change first, and it is worth seeing why. Clocked at \code{50 MHz}, a 4-bit
counter runs through all sixteen states and wraps every \code{320 ns}, roughly three million times
a second. An LED driven straight from it is not blinking, it is simply on at half brightness. To
watch a wrap you have to look at a much slower bit: count in a wider register, say 28 bits, and
assign the \textbf{top} four bits to the LEDs.

The demonstration is therefore a two-module wrapper rather than the exercise entity on its own: a
\code{reset_sync} taking the raw \code{reset_n} from the pin, and the widened counter taking its
\code{reset_s2_n} from that. The counter's port is \code{reset_s2_n} precisely so that this wiring
is the only one that type-checks, and \secref{c6:a7} repeats the pattern with a third module.
\begin{itemize}
  \item \code{clock} to the 50\,MHz oscillator, \code{reset_n} to a push button (into the
    \code{reset_sync}, never straight into the counter), and \code{count[27]} down to
    \code{count[24]} --- the top four bits of the bus Quartus builds for the widened integer port,
    since \code{count} is a \code{natural} and not a vector you could slice in VHDL --- to four
    LEDs.
  \item Bit 24 changes once every \code{2^24} clock cycles, about a third of a second, so the four
    LEDs count visibly from \code{0000} up to \code{1111} and then, with nothing in the circuit
    doing it, back to \code{0000}. A full lap takes about five seconds.
\end{itemize}
Dividing a clock by watching a high bit is the crudest possible timer, and its limitation is worth
naming now: the period is fixed at a power of two, and the only way to change it is to pick a
different bit. Chapter~\ref{c7:ch} replaces it with a counter compared against a target you choose,
which is all a timer really is.

\section{FPGA Demonstration: Shifting a Byte In by Hand}
\label{c6:a7}
\Secref{c6:a5} said a byte arrives too fast to watch. The fix is not to slow the clock down, which
would be the instinct from software and is the wrong move in hardware: the clock is the one thing
in a synchronous design you do not fiddle with. \code{serial_rx8} already has the right control,
and it is \code{shift_enable}, which ``gates the counter, not just the shift''. Hold it low and the
receiver stops dead, mid-byte, for as long as you like. Raise it for exactly one cycle and
precisely one bit goes in.

So drive \code{shift_enable} from a \textbf{button} instead of holding it high. One press, one bit.

The demonstration is a small wrapper around three modules you already have, and no new logic
beyond a single flip-flop. The course builds it live:
\begin{itemize}
  \item \code{reset_sync} (\secref{c4:a5}), turning the raw \code{reset_n} into
    \code{reset_s2_n} for everything else.
  \item \code{button_sync} (\secref{c4:a6}) at \code{COUNT = 1}, turning a press into the
    one-cycle pulse that \code{shift_enable} wants. This is not optional decoration: a raw button
    wired to \code{shift_enable} would shift a random handful of bits per press, once per clock
    cycle for as long as the contacts bounce.
  \item \code{serial_rx8} itself, unchanged, with its \code{reset_s2_n} taken from the
    \code{reset_sync} above rather than from the button. That is the whole reason the port is named
    \code{reset_s2_n}: the wrapper owns the raw \code{reset_n}, and the receiver only ever sees the
    synchronized one.
  \item One flip-flop latching \code{data_ready}, set when the pulse arrives and cleared by reset.
    Without it there is nothing to see: \code{data_ready} is high for 20\,ns, and an LED cannot
    show you that.
\end{itemize}
On the board:
\begin{itemize}
  \item \code{clock} to the 50\,MHz oscillator, \code{reset_n} to a push button.
  \item \code{serial_in} to one slide switch: this is the bit you are about to send.
  \item The shift button to a second push button.
  \item \code{data_out(7 downto 0)} to eight LEDs, and the latched \code{data_ready} to a ninth.
\end{itemize}
Now set the switch, press the button, and watch one bit enter at the right-hand LED and the whole
pattern step left. Send \code{10110010} one bit at a time and the byte assembles across the LEDs in
front of you, most significant bit first, exactly as \secref{c6:a3}'s idiom says it must. On the
eighth press the ninth LED lights and the bit counter starts over.

Three things this makes visible that the testbench can only assert:
\begin{itemize}
  \item \textbf{The bit ordering.} The first bit you send ends up in bit 7. Reading that off a
    waveform is a chore; watching it walk across a row of LEDs is not.
  \item \textbf{What \code{shift_enable} is for.} Let go of the button for a minute in the middle
    of a byte and nothing moves, nothing is lost, and the eighth bit still completes the byte. That
    is the ``pauses mid-byte rather than corrupting it'' property, demonstrated rather than claimed.
  \item \textbf{Pulse versus level, one more time.} The reason a flip-flop had to be added for
    \code{data_ready} is the whole distinction, and it is the same reason Chapter~\ref{c3:ch}'s
    edge detector existed.
\end{itemize}

% Section 6.8, from lectures/L06/README.md: the objectives, the questions to test yourself, and the
% next lecture.

\needspace{10\baselineskip}
\section{Review}
\label{c6:sec:review}

\subsection*{What you should be able to do now}
After this chapter, you should be able to:
\begin{itemize}
  \item Build a 4-bit counter as a gate network, and explain, pointing at your own drawing, why
    nothing in it makes the count return to zero.
  \item Implement the same counter in VHDL, explain wraparound as a consequence of the register's
    bit width, and say why a simulator objects to a wrap the hardware performs happily.
  \item Implement a shift register in both SIPO and PISO configurations, know which fits which
    job, and know which direction this book shifts in and why sticking to one matters.
  \item Combine a counter with a shift register to detect that a complete word has arrived, and
    verify a clocked design by running its testbench.
\end{itemize}

\needspace{6\baselineskip}
\subsection*{Questions to test yourself}
You should be able to explain:
\begin{enumerate}
  \item Why does a \code{natural range 0 to 15} counter return to \code{0} after \code{15} with no
    explicit ``if at maximum, clear it'' logic? What has to be true of the range? Which element
    performs the wrap?
  \item A VHDL simulator aborts on the wrap your CircuitVerse drawing performs without complaint.
    Which of the two is modelling the hardware, and what does that tell you about trusting either
    alone?
  \item What is the practical difference between a SIPO and a PISO shift register? Which would you
    reach for to receive a stream of serial bits, and which to drive a chain of LEDs from a fixed
    pattern?
  \item A shift register and a counter are both ``N flip-flops sharing a clock''. What differs is
    what feeds each flip-flop's input. Describe that difference for each.
  \item In \code{serial_rx8}, why is \code{data_ready} a one-cycle pulse rather than a level that
    stays high once a byte has arrived?
  \item Why does \code{shift_enable} have to gate the bit counter as well as the shift itself?
  \item The board demonstration drives \code{shift_enable} from a button rather than slowing the
    clock down. Why is slowing the clock the wrong answer, and what would break if the button were
    wired to \code{shift_enable} without passing through \code{button_sync} first?
  \item \code{data_ready} needs a flip-flop of its own before an LED can show it. Which earlier
    chapter's problem is that, in different clothes?
\end{enumerate}

\needspace{6\baselineskip}
\subsection*{Looking ahead}
Chapter~\ref{c7:ch} is about timers:
\begin{itemize}
  \item Timers: a counter with a target count, which is what turns clock ticks into real time.
  \item Built by adding exactly two things to the counter you drew here, a comparator and a clear,
    so keep that CircuitVerse project.
  \item The walking LED: this chapter's shift register, paced by that timer and started by a
    synchronized button, built live in the course and run on the DE0-CV board.
\end{itemize}

% Section 6.9, from lectures/L06/appendix/b_exercises.md.

\section{Exercises}
\label{c6:sec:exercises}

\begin{aside}
\textbf{How to work these.} Each exercise is tagged with the kinds of work it asks for:
\kinds{Paper} by hand, \kinds{CircuitVerse} in the browser, \kinds{VHDL} in a module of your own.

\textbf{How to check your work.} Every exercise below that asks you to write a VHDL module ships a
self-checking testbench under \file{lectures/L06/exercises} in the companion repository. Write your
module in its \file{exercises/<module>} directory, using the entity name and the \textbf{port
order} the exercise specifies, then run it with GHDL; \secref{c2:app:c} has the three commands.

No FPGA board is needed for any exercise. The Quartus synthesis and board-programming steps are
demonstrated in the course, and Appendix~\ref{app:quartus} shows them for a reader with a board;
your job is to get the VHDL right, and the testbench is how you confirm it. Short answers to the
parts you work on paper are in Appendix~\ref{app:answers}; reference solutions to the VHDL are
published in the companion repository, under \file{lectures/L06/solutions}. Try each exercise
yourself before looking at either.
\end{aside}

% ------------------------------------------------------------------------------------------
\exerciseset{Counters}
\label{c6:set:1}

\exercise{Sizing a counter register}{Paper}
\label{c6:ex:1}
A signal is declared as:

\begin{vhdlcode}
signal counter: natural range 0 to 1023;
\end{vhdlcode}

\begin{parts}
  \item How many bits wide is the register this synthesizes to?
  \item What is the largest value \code{counter} can hold before it wraps back around to \code{0}
    on the next increment?
  \item Rewrite the declaration for a counter that instead needs to count from \code{0} up to
    \code{63}.
\end{parts}
Follow the counter pattern in \secref{c6:a1}.

\exercise{A counter traced by hand}{Paper}
\label{c6:ex:2}
Trace a \code{natural range 0 to 3} counter by hand (no VHDL needed).

Assume:
\begin{itemize}
  \item The counter starts at \code{2}.
  \item It is incremented once per clock edge, for five consecutive edges.
  \item There is no reset in between.
\end{itemize}
Record, clock edge by clock edge:
\begin{itemize}
  \item The value held before the edge.
  \item The value captured on the edge.
\end{itemize}
Explain why the counter returns to \code{0} without any explicit ``if the count reaches its
maximum, clear it'' logic.

\note this is the \emph{synthesized} behaviour, the one that matters for hardware. \Secref{c6:a1}
flags the simulation caveat that goes with it: GHDL range-checks a \code{natural range 0 to 3}
signal, so a simulation of this counter aborts on the wrap rather than rolling over. Trace it on
paper as hardware, and reach for a full-width \code{unsigned} if you ever need a counter that wraps
in simulation too.

\exercise{Overflow in CircuitVerse}{CircuitVerse, Paper}
\label{c6:ex:3}
Build the 4-bit counter from \secref{c6:a1} by hand in CircuitVerse, and use it to see overflow
rather than take it on trust.

\begin{parts}
  \item Build it:
    \begin{itemize}
      \item A 4-bit register: four D flip-flops sharing one clock.
      \item An adder, with the register's outputs on one input and a constant \code{1} on the
        other.
      \item The adder's sum wired back to the register's \code{D} inputs.
      \item Four output elements or LEDs, one per bit, so you can read the count.
    \end{itemize}
  \item Run it at a clock period of \code{1000 ms} and record what you see:
    \begin{itemize}
      \item Write down the sequence of counts over sixteen consecutive clock edges, starting from
        \code{0000}.
      \item Watch \code{carry_out} on the top adder as the count rolls from \code{1111} to
        \code{0000}:
        \begin{itemize}
          \item On how many of the sixteen edges is it high?
          \item What is connected to it?
        \end{itemize}
      \item Explain, in one sentence, where that fifth bit goes, and why the register stores
        \code{0000} rather than \code{10000}.
    \end{itemize}
  \item Now reason about what you did \emph{not} build:
    \begin{itemize}
      \item Point at the part of your drawing that decides the counter should return to
        \code{0000}.
      \item There isn't one. Say what is doing the work instead.
      \item What would you have to add to make it count \code{0} to \code{9} and then wrap?
        (Exercise~\ref{c6:ex:4} asks you to write exactly that in VHDL, so keep your answer.)
    \end{itemize}
  \item Change the width to three bits by removing the top flip-flop and narrowing the adder:
    \begin{itemize}
      \item Predict the new wrap point before you run it.
      \item Confirm it.
      \item State the rule connecting a counter's width to the value it wraps at.
    \end{itemize}
\end{parts}
\textbf{Save this circuit.} Chapter~\ref{c7:ch} builds its timer by adding two things to it --- a
comparator and a clear --- and starting from a counter you have already watched work is the whole
point of the exercise there.

\exercise{A counter of any radix}{VHDL}
\label{c6:ex:4}
Write an entity named \code{counter} that counts from \code{0} up to one below its radix, then wraps
back to \code{0}. With the default radix that is a decimal counter, \code{0} through \code{9}.

The entity has one generic:

\begin{booktable}{@{}p{3.6em}p{9.2em}p{3.9em}L@{}}
  \thead{Generic} & \thead{Type} & \thead{Default} & \thead{Description} \\
  \midrule
  \code{RADIX} & \code{natural range 1 to 2**16} & \code{10} & How many values the counter passes
    through before repeating. It counts \code{0} to \code{RADIX-1} inclusive, so \code{RADIX} is
    the length of the cycle, not the largest count. \\
\end{booktable}

\noindent and these ports:

\begin{booktable}{@{}p{5.6em}p{4.8em}p{9.2em}L@{}}
  \thead{Port} & \thead{Direction} & \thead{Type} & \thead{Description} \\
  \midrule
  \code{clock} & in & \code{std_logic} & System clock. \\
  \code{reset_s2_n} & in & \code{std_logic} & Active-low, \textbf{already-synchronized}: it comes
    from your \code{reset_sync} (Chapter~4's Exercise~\ref{c4:ex:8}), never straight from a pin.
    Asynchronous, so it clears with no clock edge involved. \\
  \code{count} & out & \code{natural range 0 to RADIX-1} & The running count. \code{0} is a value
    it takes, so the range starts there: reset drives it to \code{0}, and it returns there on
    every wrap. \\
  \code{tick} & out & \code{std_logic} & A one-cycle pulse on the edge where \code{count} wraps
    from \code{RADIX-1} back to \code{0}. \\
\end{booktable}

\noindent Implement it using a single synchronous process:
\begin{itemize}
  \item Include \code{clock} and \code{reset_s2_n} in the sensitivity list.
  \item On reset, set \code{count} to \code{0} and \code{tick} to \code{'0'}.
  \item On each rising clock edge:
    \begin{itemize}
      \item Drive \code{tick} low by default.
      \item When \code{count} has reached \code{RADIX-1}, wrap \code{count} to \code{0} and pulse
        \code{tick} high.
      \item Otherwise, increment \code{count}.
    \end{itemize}
\end{itemize}
Write the comparison against \code{RADIX-1}, never against \code{9}. The generic is the whole
difference between a module and a module you can use twice: \code{RADIX} is fixed before synthesis
runs, so the comparator it builds is no larger than a hard-coded one, and the counter's register is
sized to the radix you ask for. This is the same argument \secref{c4:a8} makes for
\code{button_sync}'s \code{COUNT}, and the reason Chapter~\ref{c7:ch}'s \code{timer} takes its
target as a generic rather than as a constant.

Unlike the free wraparound of a \code{natural range 0 to 15} counter (\secref{c6:a1}), a counter of
arbitrary radix needs an explicit comparison, because \code{RADIX-1} is not in general a
power-of-two boundary. Notice what happens when it is: with \code{RADIX = 16} the comparison is
still written, but the hardware underneath no longer needs it.

\note[Tip] The process needs to read the running count (to compare it and to increment it), so hold
the count in an internal \code{signal count_s: natural range 0 to RADIX-1;}, run the logic on
\code{count_s}, and drive the output with \code{count <= count_s;} outside the process, exactly as
\file{d_flip_flop.vhd} uses \code{q_s} in Chapter~\ref{c3:ch}. This is not a style choice: every
\code{ghdl} command in this book passes \code{--std=93}, and under VHDL-93 reading an \code{out}
port does not compile at all. Later standards relax that, which is why you may have seen it done
elsewhere.

Follow the counter pattern in \secref{c6:a1}, and the drive-low-then-pulse shape of
\code{data_ready} in \secref{c6:a5}.

\modulefigure{counter}

\selfcheck name your entity \code{counter}, with generic \code{RADIX} (\code{natural range 1 to
2**16}, default \code{10}), inputs \code{clock}, \code{reset_s2_n} and outputs \code{count}
(\code{natural range 0 to RADIX-1}) and \code{tick} (\code{std_logic}), declared in that order; its
testbench is in \file{exercises/counter}. It elaborates \textbf{two} instances, one at
\code{RADIX = 10} and one at \code{RADIX = 4}, and runs each through two full cycles, checking that
\code{tick} is a one-cycle pulse rather than a level and that it never fires part way through. An
architecture that quietly hard-codes \code{10} passes the first instance and fails the second,
which is the point of checking two.

% ------------------------------------------------------------------------------------------
\exerciseset{Shift Registers}
\label{c6:set:2}

\exercise{Tracing a SIPO register}{Paper}
\label{c6:ex:5}
An 8-bit SIPO shift register starts at \code{00000000} and uses the idiom from \secref{c6:a3}:

\begin{vhdlcode}
shift_reg <= shift_reg(6 downto 0) & serial_in;
\end{vhdlcode}

\noindent It receives the serial bit stream \code{1, 0, 1, 1} (one bit per clock edge, in that
order) over four consecutive rising clock edges.

\begin{parts}
  \item Write out the full 8-bit contents of \code{shift_reg} after each of the four clock edges.
  \item After those four edges:
    \begin{itemize}
      \item Which four bits of \code{shift_reg} hold the received data?
      \item In what order do they appear, relative to the order the bits arrived in?
    \end{itemize}
\end{parts}

\exercise{Tracing a PISO register}{Paper}
\label{c6:ex:6}
A 4-bit PISO register, following the idiom from \secref{c6:a4}, is:
\begin{itemize}
  \item Parallel-loaded with \code{parallel_in = "1010"} (\code{load = '1'} for one clock edge).
  \item Then shifted with \code{shift = '1'} for the next four clock edges.
  \item With \code{serial_in} held at \code{'0'} throughout.
\end{itemize}
Remember the direction from \secref{c6:a2}:
\begin{itemize}
  \item bits move toward the MSB.
  \item \code{serial_out} is bit 3, the value about to leave.
  \item \code{serial_in} enters at bit 0.
\end{itemize}

\begin{parts}
  \item Write out the register's contents and \code{serial_out} at each step, one line per clock
    edge:
    \begin{itemize}
      \item The line for the load edge, before any shifting.
      \item Then one line for each of the four shift edges.
    \end{itemize}
  \item In what order do the four loaded bits appear on \code{serial_out}, relative to their
    positions in \code{parallel_in}?
  \item What are the final contents of the register after all four shifts, and why?
\end{parts}
\note[Tip] You don't need real hardware to check Exercises~\ref{c6:ex:5} and~\ref{c6:ex:6} by hand.
Walk through them exactly as you would trace a truth table: one clock edge, one line, at a time.

\exercise{SIPO and PISO in VHDL}{VHDL, Paper}
\label{c6:ex:7}
Implement the two useful shift-register configurations from this chapter in VHDL.

\begin{parts}
  \item Write an entity named \code{sipo8}, a serial-in/parallel-out register.

    The entity must have:

    \begin{booktable}{@{}p{6.3em}p{4.8em}p{9.2em}L@{}}
      \thead{Port} & \thead{Direction} & \thead{Type} & \thead{Description} \\
      \midrule
      \code{clock} & in & \code{std_logic} & System clock. \\
      \code{reset_s2_n} & in & \code{std_logic} & Active-low, \textbf{already-synchronized}: it
        comes from your \code{reset_sync} (Chapter~4's Exercise~\ref{c4:ex:8}), never straight
        from a pin. Asynchronous, so it clears with no clock edge involved. Clears the whole
        register, so \code{parallel_out} reads \code{"00000000"} while it is held low. \\
      \code{serial_in} & in & \code{std_logic} & The incoming bit. \\
      \code{parallel_out} & out & \code{std_logic_vector(7 downto 0)} & The register's
        contents. \\
    \end{booktable}

    Follow the SIPO idiom from \secref{c6:a3}.
\end{parts}

\modulefigure{sipo8}

\selfcheck name your entity \code{sipo8}, with ports \code{clock}, \code{reset_s2_n},
\code{serial_in} (in) and \code{parallel_out} (\code{std_logic_vector(7 downto 0)}, out), declared
in that order; its testbench is in \file{exercises/sipo8}.

\begin{parts}\setcounter{enumi}{1}
  \item Write an entity named \code{piso8}, a parallel-in/serial-out register.

    The entity must have:

    \begin{booktable}{@{}p{6.3em}p{4.8em}p{9.2em}L@{}}
      \thead{Port} & \thead{Direction} & \thead{Type} & \thead{Description} \\
      \midrule
      \code{clock} & in & \code{std_logic} & System clock. \\
      \code{reset_s2_n} & in & \code{std_logic} & Active-low, \textbf{already-synchronized}: it
        comes from your \code{reset_sync} (Chapter~4's Exercise~\ref{c4:ex:8}), never straight
        from a pin. Asynchronous, so it clears with no clock edge involved. Clears the whole
        register, so \code{serial_out} reads \code{'0'} while it is held low. \\
      \code{load} & in & \code{std_logic} & Capture \code{parallel_in} on the next rising
        edge. \\
      \code{shift} & in & \code{std_logic} & Shift one position on the next rising edge. \\
      \code{serial_in} & in & \code{std_logic} & The bit fed in at bit 0 on a shift. \\
      \code{parallel_in} & in & \code{std_logic_vector(7 downto 0)} & The value captured on a
        load. \\
      \code{serial_out} & out & \code{std_logic} & The register's top bit. \\
    \end{booktable}

    Follow the PISO idiom from \secref{c6:a4}:
    \begin{itemize}
      \item On a load edge (\code{load = '1'}), capture \code{parallel_in}.
      \item On a shift edge (\code{shift = '1'}), move every bit one position toward the MSB,
        feeding \code{serial_in} in at bit 0.
      \item Drive \code{serial_out} from bit 7, so a loaded byte leaves most-significant bit first.
      \item \code{load} takes priority over \code{shift} when both are high.
      \item With both low, the register \textbf{holds}: nothing moves and nothing is captured. Its
        testbench checks this, so a design that shifts unconditionally fails.
    \end{itemize}
\end{parts}

\modulefigure{piso8}

\selfcheck name your entity \code{piso8}, with ports \code{clock}, \code{reset_s2_n}, \code{load},
\code{shift}, \code{serial_in}, \code{parallel_in} (\code{std_logic_vector(7 downto 0)}) and
\code{serial_out} (out), declared in that order; its testbench is in \file{exercises/piso8}. Keep
this module: Chapter~\ref{c8:ch}'s transmitter capstone (Exercise~\ref{c8:ex:7}) does not
instantiate it, but it is worth rereading there for the idiom, since that machine shifts the byte
out inline.

Re-derive each idiom from the chapter rather than copying \code{serial_rx8}'s shift line.

\begin{parts}\setcounter{enumi}{2}
  \item For each of the four shift-register configurations, name one task it is the natural choice
    for:
    \begin{itemize}
      \item SISO (serial-in/serial-out).
      \item SIPO (serial-in/parallel-out).
      \item PISO (parallel-in/serial-out).
      \item PIPO (parallel-in/parallel-out).
    \end{itemize}
    Follow the summary table in \secref{c6:a2}.
\end{parts}

% ------------------------------------------------------------------------------------------
\exerciseset{Reading the Serial Receiver}
\label{c6:set:3}

The worked example, \code{serial_rx8} (\file{lectures/L06/serial_rx8/serial_rx8.vhd}), ships a
reference testbench. Run it from the example's own directory before starting:

\begin{shell}
cd lectures/L06/serial_rx8
ghdl -a --std=93 serial_rx8.vhd serial_rx8_tb.vhd
ghdl -e --std=93 serial_rx8_tb
ghdl -r --std=93 serial_rx8_tb --assert-level=error --stop-time=10ms
\end{shell}

\exercise{Accounting for the receiver}{Paper, VHDL}
\label{c6:ex:8}
Account for the receiver's behaviour, referring to \secref{c6:a5}.

\begin{parts}
  \item The bit stream \code{1, 0, 1, 1, 0, 0, 1, 0} arrives, one bit per enabled clock edge:
    \begin{itemize}
      \item What is on \code{data_out} in the cycle where \code{data_ready} pulses?
      \item Which arriving bit ends up in \code{data_out(7)}, and which in \code{data_out(0)}?
    \end{itemize}
  \item \code{shift_enable} goes low for three clock cycles after the fourth bit, then high again
    for the remaining four:
    \begin{itemize}
      \item Does the receiver still produce the same byte? Explain what the bit counter does during
        the gap.
      \item What would go wrong instead if \code{shift_enable} gated only the shift register and
        not the counter?
    \end{itemize}
  \item Make each of the following changes to your own copy of \file{serial_rx8.vhd}, one at a
    time, and \textbf{predict} what the reference testbench will report before you run it. Then run
    it and account for the result:
    \begin{itemize}
      \item Change the shift line to the mirror-image direction,
        \code{shift_reg <= serial_in & shift_reg(7 downto 1);}.
      \item Remove the unconditional \code{data_ready <= '0';} from the top of the clocked branch.
      \item Change the bit-counter comparison from \code{BIT_WIDTH - 1} to \code{BIT_WIDTH - 2},
        that is from 7 to 6.
    \end{itemize}
    Each of these is a real bug someone has shipped. For each one, say which of the testbench's
    checks catches it, and what the failure would have looked like on a real serial link if nothing
    had caught it.
\end{parts}

